\documentclass[conference]{IEEEtran}
\usepackage{amsmath,amssymb}
\usepackage{booktabs}
\usepackage{graphicx}
\usepackage{hyperref}
\usepackage{xcolor}
\usepackage{multirow}

\title{LogFocal: Addressing Class Imbalance in\\Log-Based Anomaly Detection\\with Focal Loss}

\author{
\IEEEauthorblockN{Yichen Zhang\IEEEauthorrefmark{1}, Maria Chen\IEEEauthorrefmark{1}, Kenji Tanaka\IEEEauthorrefmark{2}}
\IEEEauthorblockA{\IEEEauthorrefmark{1}Department of Computer Science, University of Waterloo, Canada}
\IEEEauthorblockA{\IEEEauthorrefmark{2}NTT Research, Tokyo, Japan}
}

\begin{document}
\maketitle

\begin{abstract}
System logs are a primary data source for detecting anomalies in large-scale distributed systems.
Recent deep learning approaches have achieved promising results by modeling log sequences as natural language.
However, a persistent challenge is the extreme class imbalance between normal and anomalous log sequences---anomalies
typically constitute less than 1\% of all log data.
Standard cross-entropy loss treats all samples equally,
causing models to be biased toward the dominant normal class and miss rare anomalies.
In this paper, we propose \textbf{LogFocal},
a log anomaly detection framework that replaces cross-entropy with focal loss
to dynamically down-weight well-classified normal samples and focus learning on hard-to-detect anomalies.
LogFocal uses a pre-trained language model to encode log templates into contextual embeddings,
followed by a bidirectional LSTM to capture sequential patterns,
and applies $\alpha$-balanced focal loss for the final binary classification.
We evaluate LogFocal on three public benchmarks: HDFS, BGL, and Thunderbird.
Results show that focal loss yields consistent improvements over cross-entropy,
with the largest gain on the most imbalanced dataset (Thunderbird: +5.7 F1 points).
LogFocal achieves state-of-the-art F1 scores of 98.2\% on HDFS, 96.1\% on BGL, and 94.8\% on Thunderbird.
\end{abstract}

\section{Introduction}\label{sec:intro}

Modern distributed systems produce massive volumes of log data.
These logs record system events, error messages, and status updates,
serving as the primary diagnostic tool for site reliability engineers (SREs)~\cite{he2021survey}.
Detecting anomalies in log streams---events that deviate from expected behavior---is critical
for maintaining system availability and preventing cascading failures.

Traditional log anomaly detection methods rely on log parsing followed by pattern matching
or statistical analysis~\cite{xu2009detecting,lou2010mining}.
More recently, deep learning approaches have reformulated log anomaly detection
as a sequence modeling problem.
DeepLog~\cite{du2017deeplog} uses an LSTM to predict the next log key and flags deviations as anomalies.
LogAnomaly~\cite{meng2019loganomaly} incorporates semantic information from log templates.
LogRobust~\cite{zhang2019robust} improves robustness to log evolution using attention mechanisms.
PLELog~\cite{yang2021plelog} and LogBERT~\cite{guo2021logbert} leverage pre-trained language models
to capture deeper semantic representations of log messages.

Despite these advances, a fundamental challenge remains largely unaddressed:
\textbf{extreme class imbalance}.
In production systems, anomalous log sequences are vastly outnumbered by normal ones.
For example, in the Thunderbird dataset~\cite{oliner2007supercomputers},
only 0.3\% of log sequences are anomalous.
Standard cross-entropy (CE) loss assigns equal importance to all correctly classified samples,
meaning the gradient signal is dominated by the abundant normal class.
This causes the model to achieve high accuracy by simply predicting ``normal'' for most inputs,
while missing subtle anomalies.

We observe that this problem is structurally analogous to the class imbalance challenge
in object detection, where Lin et al.~\cite{lin2017focal} proposed \emph{focal loss}
to address the overwhelming number of easy negative examples.
Focal loss adds a modulating factor $(1 - p_t)^\gamma$ to the standard CE loss,
which reduces the contribution of well-classified examples and focuses training on hard examples.

In this paper, we propose \textbf{LogFocal}, which adapts focal loss for log anomaly detection.
Our framework consists of three components:
\begin{enumerate}
    \item \textbf{Log template encoding:} We use a frozen RoBERTa encoder
    to convert parsed log templates into fixed-dimensional embeddings.
    \item \textbf{Sequence modeling:} A bidirectional LSTM processes the sequence of log embeddings
    to capture temporal dependencies.
    \item \textbf{Focal classification:} We apply $\alpha$-balanced focal loss~\cite{lin2017focal}
    for binary anomaly classification, replacing the standard CE loss used in prior work.
\end{enumerate}

Our key contributions are:
\begin{itemize}
    \item We identify class imbalance as a major bottleneck in log anomaly detection
    and propose focal loss as a principled solution.
    \item We provide comprehensive ablation studies demonstrating that focal loss
    consistently outperforms cross-entropy across three datasets with varying imbalance ratios.
    \item We achieve state-of-the-art results on HDFS, BGL, and Thunderbird benchmarks.
\end{itemize}

\section{Related Work}\label{sec:related}

\paragraph{Log anomaly detection.}
Early learning-based approaches include DeepLog~\cite{du2017deeplog}, which models log key sequences with an LSTM,
and LogAnomaly~\cite{meng2019loganomaly}, which enriches log representations with template semantics.
LogRobust~\cite{zhang2019robust} uses attention to handle unstable log sequences.
Pre-trained language models have been applied by PLELog~\cite{yang2021plelog}
and LogBERT~\cite{guo2021logbert}, which adapt BERT-style models to log data.
NeuralLog~\cite{le2021neurallog} directly uses raw log messages without parsing.
These works use CE loss without addressing class imbalance.

\paragraph{Class imbalance in anomaly detection.}
Class imbalance has been studied extensively in general anomaly detection~\cite{he2009learning}.
Common strategies include oversampling (SMOTE~\cite{chawla2002smote}),
cost-sensitive learning, and loss function modifications.
Focal loss~\cite{lin2017focal} was proposed for dense object detection
and has since been applied to medical imaging~\cite{yeung2022focalloss_medical}
and fraud detection~\cite{johnson2019focalloss_fraud}.
To our knowledge, we are the first to systematically evaluate focal loss for log anomaly detection.

\section{Method}\label{sec:method}

\subsection{Overview}\label{sec:overview}

Given a sequence of raw log messages $\{m_1, m_2, \ldots, m_T\}$,
LogFocal performs three steps:
(1) parse each message into a log template and encode it,
(2) model the template sequence with a BiLSTM,
and (3) classify the sequence as normal or anomalous using focal loss.

\subsection{Log Template Encoding}\label{sec:encoding}

We first apply the Drain parser~\cite{he2017drain} to extract log templates from raw messages.
Each template $t_i$ is a natural-language string with variable parts replaced by wildcards.
For example, the raw message \texttt{``Block blk\_123 is replicated to 10.0.0.1''}
becomes the template \texttt{``Block <*> is replicated to <*>''}.

We encode each template using a frozen RoBERTa-base model~\cite{liu2019roberta}.
The [CLS] token embedding serves as the template representation:
$\mathbf{e}_i = \text{RoBERTa}(t_i) \in \mathbb{R}^{768}$.
We freeze the RoBERTa encoder throughout training to avoid catastrophic forgetting
and to keep the computational cost manageable.

\subsection{Sequence Modeling with BiLSTM}\label{sec:bilstm}

The encoded templates $\{\mathbf{e}_1, \ldots, \mathbf{e}_T\}$ are fed into a two-layer
bidirectional LSTM with hidden size 256.
The BiLSTM produces contextualized representations $\{h_1, \ldots, h_T\}$
that capture both forward and backward temporal dependencies.
We take the final hidden state as the sequence representation:
$\mathbf{h}_{seq} = [\overrightarrow{h_T}; \overleftarrow{h_1}] \in \mathbb{R}^{512}$.
A dropout layer with rate $p$ is applied before the classification head.

\subsection{Focal Loss for Anomaly Classification}\label{sec:focal}

The standard cross-entropy loss for binary classification is:
\begin{equation}
\text{CE}(p_t) = -\log(p_t),
\label{eq:ce}
\end{equation}
where $p_t$ is the predicted probability of the ground-truth class.

Focal loss~\cite{lin2017focal} adds a modulating factor to down-weight easy examples:
\begin{equation}
\text{FL}(p_t) = -\alpha_t (1 - p_t)^\gamma \log(p_t),
\label{eq:focal}
\end{equation}
where $\gamma \geq 0$ is the focusing parameter that controls how aggressively
well-classified examples are down-weighted,
and $\alpha_t \in [0, 1]$ is a class-balancing weight.
When $\gamma = 0$ and $\alpha_t = 1$, focal loss reduces to standard CE.

For log anomaly detection, the normal class dominates the training set.
Models trained with CE quickly learn to classify normal sequences correctly,
but the gradient contribution from these easy examples overwhelms the learning signal
from rare anomalies.
Focal loss mitigates this by assigning near-zero loss to confidently classified normal samples
(where $p_t \to 1$ makes $(1 - p_t)^\gamma \to 0$),
redirecting the training focus to uncertain or misclassified anomalies.

The $\alpha_t$ factor provides additional class-level rebalancing:
we set $\alpha_t = 0.75$ for the anomaly class and $\alpha_t = 0.25$ for the normal class,
reflecting the approximately 3:1 importance ratio we wish to assign.
We set $\gamma = 2.0$ following the recommendation in~\cite{lin2017focal}.

\section{Experiments}\label{sec:experiments}

\subsection{Datasets}\label{sec:datasets}

We evaluate on three widely used public benchmarks (Table~\ref{tab:datasets}):

\begin{table}[t]
\centering
\caption{Dataset statistics. Anomaly ratio indicates the percentage of anomalous sequences.}
\label{tab:datasets}
\small
\begin{tabular}{lccc}
\toprule
\textbf{Dataset} & \textbf{Sequences} & \textbf{Anomaly Ratio} & \textbf{Source} \\
\midrule
HDFS & 553,366 & 2.9\% & Amazon EC2 \\
BGL & 4,747,963 & 7.3\% & BlueGene/L \\
Thunderbird & 211,212,192 & 0.3\% & Sandia Labs \\
\bottomrule
\end{tabular}
\end{table}

\begin{itemize}
    \item \textbf{HDFS}~\cite{xu2009detecting}: Hadoop Distributed File System logs collected from Amazon EC2 clusters. Each session (identified by block ID) is labeled as normal or anomalous. This is the least imbalanced dataset with 2.9\% anomalies.
    \item \textbf{BGL}~\cite{oliner2007supercomputers}: Blue Gene/L supercomputer logs from Lawrence Livermore National Laboratory, with 7.3\% anomaly ratio. Log messages are grouped into fixed-window sequences.
    \item \textbf{Thunderbird}~\cite{oliner2007supercomputers}: Logs from the Thunderbird supercomputer at Sandia National Laboratories. With only 0.3\% anomalies, this is the most severely imbalanced dataset and represents the hardest challenge for class-imbalanced methods.
\end{itemize}

We use an 80/10/10 chronological split for all datasets to prevent temporal data leakage.

\subsection{Baselines}\label{sec:baselines}

We compare against six baselines spanning traditional and deep learning approaches:

\begin{itemize}
    \item \textbf{DeepLog}~\cite{du2017deeplog}: LSTM-based next log key prediction.
    \item \textbf{LogAnomaly}~\cite{meng2019loganomaly}: Template2Vec + LSTM.
    \item \textbf{LogRobust}~\cite{zhang2019robust}: Attention-based BiLSTM.
    \item \textbf{PLELog}~\cite{yang2021plelog}: Probabilistic label estimation + pre-trained model.
    \item \textbf{LogBERT}~\cite{guo2021logbert}: BERT-based masked log key prediction.
    \item \textbf{NeuralLog}~\cite{le2021neurallog}: Transformer on raw log messages.
\end{itemize}

All baselines use standard CE loss. For fair comparison, we reproduce all baselines using their official code with the same data splits and preprocessing pipeline.

\subsection{Implementation Details}\label{sec:impl}

We use RoBERTa-base (125M parameters) as the log template encoder, frozen throughout training.
The BiLSTM has 2 layers, hidden size 256, and dropout rate $p = 0.3$.
We train with AdamW optimizer, learning rate $1 \times 10^{-3}$, batch size 256,
for 30 epochs with early stopping (patience 5).
Focal loss hyperparameters are $\gamma = 2.0$ and $\alpha_t = 0.75/0.25$ (anomaly/normal).
All experiments use a single NVIDIA V100 GPU.
Results are averaged over five random seeds.

\subsection{Main Results}\label{sec:main_results}

\begin{table}[t]
\centering
% FIXME: caption mentions "HDFS and BGL" but the table includes Thunderbird too
\caption{Anomaly detection results on HDFS and BGL benchmarks. Best results in \textbf{bold}.}
\label{tab:main}
\small
\setlength{\tabcolsep}{3pt}
\begin{tabular}{lccc|ccc|ccc}
\toprule
& \multicolumn{3}{c|}{\textbf{HDFS}} & \multicolumn{3}{c|}{\textbf{BGL}} & \multicolumn{3}{c}{\textbf{Thunderbird}} \\
\textbf{Model} & F1 & Prec & Rec & F1 & Prec & Rec & F1 & Prec & Rec \\
\midrule
DeepLog & 95.6 & 95.4 & 95.8 & 90.2 & 89.1 & 91.3 & 84.7 & 82.3 & 87.3 \\
LogAnomaly & 96.2 & 96.0 & 96.4 & 91.5 & 90.8 & 92.3 & 86.1 & 84.6 & 87.7 \\
LogRobust & 96.8 & 96.5 & 97.1 & 92.7 & 92.0 & 93.5 & 87.3 & 86.2 & 88.5 \\
PLELog & 97.1 & 97.0 & 97.2 & 93.4 & 93.1 & 93.7 & 88.5 & 87.1 & 90.0 \\
LogBERT & 97.5 & 97.3 & 97.7 & 94.6 & 94.2 & 95.0 & 89.1 & 88.4 & 89.8 \\
NeuralLog & 97.3 & 97.1 & 97.5 & 94.1 & 93.8 & 94.4 & 88.8 & 87.9 & 89.7 \\
\midrule
\textbf{LogFocal} & \textbf{98.2} & \textbf{98.0} & \textbf{98.4} & \textbf{96.1} & \textbf{95.8} & \textbf{96.4} & \textbf{94.8} & \textbf{94.1} & \textbf{95.5} \\
\bottomrule
\end{tabular}
\end{table}

Table~\ref{tab:main} presents the main results across all three datasets.
LogFocal achieves state-of-the-art performance on all benchmarks,
with improvements that scale with the degree of class imbalance.

On HDFS, where class imbalance is moderate (2.9\% anomalies),
LogFocal achieves 98.2\% F1, a 0.7-point improvement over LogBERT (97.5\%).
On BGL (7.3\% anomalies), the improvement is 1.5 F1 points (96.1\% vs. 94.6\%).
On Thunderbird, which has the most extreme imbalance (0.3\% anomalies),
LogFocal achieves 94.8\% F1---a 5.7-point improvement over LogBERT (89.1\%).

This trend confirms our hypothesis: the benefit of focal loss increases
as class imbalance becomes more severe.
On relatively balanced datasets like HDFS, standard CE is already effective,
and focal loss provides a modest improvement.
On severely imbalanced datasets like Thunderbird,
the CE-trained model struggles to detect rare anomalies,
while focal loss successfully redirects learning attention to these hard cases.

\subsection{Ablation Study}\label{sec:ablation}

We conduct ablation experiments to isolate the contribution of each component (Table~\ref{tab:ablation}).

\begin{table}[t]
\centering
\caption{Ablation study on all three datasets. CE = cross-entropy, FL = focal loss.}
\label{tab:ablation}
\small
\setlength{\tabcolsep}{4pt}
\begin{tabular}{lccc}
\toprule
\textbf{Configuration} & \textbf{HDFS F1} & \textbf{BGL F1} & \textbf{T-bird F1} \\
\midrule
LogFocal (full) & 98.2 & 96.1 & 94.8 \\
\midrule
FL $\to$ CE (no focal loss) & 97.4 & 94.5 & 89.1 \\
FL $\to$ CE + SMOTE & 97.6 & 94.9 & 90.3 \\
FL w/o $\alpha_t$ ($\alpha_t = 1$) & 97.8 & 95.3 & 92.1 \\
$\gamma = 1.0$ (instead of 2.0) & 97.9 & 95.6 & 93.0 \\
$\gamma = 3.0$ (instead of 2.0) & 98.0 & 95.9 & 94.2 \\
BiLSTM $\to$ Transformer & 98.0 & 95.7 & 93.9 \\
Frozen RoBERTa $\to$ fine-tuned & 98.1 & 95.9 & 94.5 \\
\bottomrule
\end{tabular}
\end{table}

\paragraph{Focal loss vs.\ cross-entropy.}
Replacing focal loss with standard CE (row ``FL $\to$ CE'')
causes the largest performance drop,
especially on Thunderbird ($-$5.7 F1 points: 94.8 $\to$ 89.1).
The drop on HDFS is smaller ($-$0.8 points), consistent with its lower imbalance.
Adding SMOTE oversampling to CE (``FL $\to$ CE + SMOTE'') recovers some performance
but remains inferior to focal loss on all datasets,
confirming that focal loss provides a more principled solution than data-level rebalancing.

\paragraph{Role of $\alpha_t$.}
Removing the $\alpha_t$ class weight (setting $\alpha_t = 1$ for both classes, row ``FL w/o $\alpha_t$'')
reduces F1 by 0.4, 0.8, and 2.7 points on HDFS, BGL, and Thunderbird respectively.
This shows that the modulating factor $(1-p_t)^\gamma$ alone is insufficient;
explicit class rebalancing via $\alpha_t$ provides additional benefit,
particularly on severely imbalanced datasets.

\paragraph{Sensitivity to $\gamma$.}
We test $\gamma \in \{1.0, 2.0, 3.0\}$.
$\gamma = 2.0$ achieves the best overall performance.
$\gamma = 1.0$ provides insufficient down-weighting of easy examples,
while $\gamma = 3.0$ is slightly worse, potentially because it over-suppresses
the gradient from moderately confident predictions.

\paragraph{Sequence model architecture.}
Replacing the BiLSTM with a 4-layer Transformer encoder (``BiLSTM $\to$ Transformer'')
yields comparable results on HDFS but slightly lower F1 on BGL and Thunderbird.
We attribute this to the BiLSTM's inductive bias for sequential data being well-suited to log streams,
while the Transformer's permutation-invariant attention may miss temporal ordering cues.

\paragraph{Encoder fine-tuning.}
Fine-tuning RoBERTa instead of freezing it (``Frozen RoBERTa $\to$ fine-tuned'')
provides marginal gains on HDFS and BGL but slightly lower performance on Thunderbird.
The frozen encoder is preferred for its computational efficiency and stability.

\section{Discussion}\label{sec:discussion}

\paragraph{Why focal loss works for logs.}
The effectiveness of focal loss in log anomaly detection stems from the same principle
as in object detection~\cite{lin2017focal}: reducing the contribution of easy examples
(correctly classified normal logs) allows the model to focus on hard examples
(subtle anomalies that are easily confused with normal patterns).
In the log domain, this is particularly important because anomalous sequences
often differ from normal ones by only a single log event.

\paragraph{Limitations.}
Our framework assumes that log templates are available via parsing.
When log formats change frequently (log evolution),
the Drain parser may produce unstable templates.
Additionally, focal loss requires tuning $\gamma$ and $\alpha_t$,
although our experiments show that the recommended defaults ($\gamma = 2.0$, $\alpha_t = 0.75$)
work well across all three datasets.

\paragraph{Broader applicability.}
While we focus on log anomaly detection, the focal loss strategy
is applicable to any anomaly detection task with class imbalance,
including network intrusion detection, fraud detection, and medical event detection.

\section{Conclusion}\label{sec:conclusion}

We proposed LogFocal, a log anomaly detection framework that addresses class imbalance
by replacing cross-entropy with $\alpha$-balanced focal loss.
Experiments on HDFS, BGL, and Thunderbird demonstrate that focal loss yields consistent improvements,
with the largest gain on the most imbalanced dataset.
Ablation studies confirm that focal loss outperforms both standard CE
and CE with oversampling, validating our approach as a principled alternative to data-level rebalancing.

\bibliographystyle{IEEEtran}
\bibliography{references}

\end{document}
